\section{MRX电子层：确认了双尺度，也留下绝对量的差异}
\citet{L119}使用五列磁探针测$B_T$，径向最高分辨2.5 mm，轴向约3 cm。
由Ampère定律反演面内电流，再以$\vel_e\simeq-\J/(en)$求电子流；Mach探针检查离子流是否可忽略。
因此电子速度是由磁场导数和密度反演的量，不能写成直接电子分布测量。
其H、D、He放电中，Hall四极幅度可达上游重联场约30\%--60\%，电子排出通道比离子窄且更快。
但只有约一半放电出现可辨四极场，文章选这些事件作分析；这限制其统计结论的适用总体。

电子层宽$\delta_{BT}$定义为速度峰到降至峰值40\%位置的径向距离，并取轴向排出速度峰处的值。
它约为$8d_e$，而比较二维模型常为$1$--$2d_e$；峰值排出约$0.11V_{\!Ae}$，$V_{\!Ae}$所用磁场在电子层边缘取值。
宽度乘速度的归一化乘积约$8\times0.11=0.88$，支持较宽较慢通道仍能输送相近电子通量的图像。
长度约$40$--$80d_e$，随密度变化的趋势不支持简单$L_e\propto d_e$，更不能据大误差数据宣布严格反比例律。

图13中，重联场约170 V/m，经典电阻约40 V/m，另一个未细测Hall分量估约50 V/m，所估对流惯性约1 V/m。
压力张量及涨落项尚未全测，欧姆预算并未闭合。
第一章摩擦模型能把电子排出拉低，但有效碰撞率来自$E/J$反推，并假设同样作用于排出方向；仍可能高估约两倍。
这项研究最有价值的未决结果就是：\emph{测到正确的电子尺度缩放，不等于已经解释绝对层厚与速度}。

对应PIC模拟纳入通量芯、粒子吸收与壁面，质量比75，并以整体率和离子层厚选择参数。
较好的流型吻合强调边界的重要性，也意味着它不是完全不依赖实验输入的预测。
流体后续验证应使用同样的诊断定义和仪器卷积，而不是用理想无限Harris盒与实验数值直接比高低。

\section{UTST合并：势降与排出同时出现，加热仍需独立检验}
\citet{L189}在球形托卡马克合并中同步比较磁探针、浮动电势径向扫描和Doppler光谱。
在约200--205微秒，重联磁通变化加快、电流降低，X点附近浮动电势梯度变陡；覆盖X点的视线看到接近极向Alfvén速度的离子流。
典型重联场$B_p=0.01$ T、导场$B_t=0.1$ T，电场估400--800 V/m，碳离子流约27 km/s，极向Alfvén速度约24 km/s。
用$\tau_{\rm acc}=m_iv_i/(q_iE_r)$估加速时间2.1--4.2微秒，与碳离子回旋周期3.9微秒可比，支持瞬态静电加速的可能性。

但浮动电势并不就是空间电势：二者相差依赖$T_e$的鞘层项，只有空间温度差较小才能把其梯度当电场。
本文未给完整同步温度修正；视线积分的谱宽也混有流动展宽，作者对此明确保留。
正文CIII与图6 HeII标注不一致，进一步限制精确的物种温度解释。
故直接证据是势降、近Alfvén流及重联时序的联系；“Hall独自导致不可逆离子加热”和具体热化机制仍是待验证解释。

\section{三维EMHD磁涡：没有观察到细电流片同样值得记录}
\citet{F766}在大尺寸氩等离子体中用天线先建立反向偶极磁场，关断电流后观察自由松弛。
三维磁场来自可重复放电中移动磁探针的扫描，空间步约0.75 cm、时间分辨10 ns，电子惯性长度约0.75 cm。
初始cusp零点在等离子体电流作用下变成不规则spiral零点；闭合磁通随后转为开放。
双磁涡还经退化零点、零线的变化合并成更长寿命FRC，这是区别于二维X点图像的重要实验现象。

Langmuir探针显示电子温度约倍增，发光主要发生在天线关断后的松弛期间；反转天线极性而不形成同类零线时，发光弱得多。
实验没有观察到强whistler湍流、显著异常电阻或细至$d_e$的电流片；由回路电压和电流估计的导电率与经典值接近。
这些负结果值得保留，不能为配合某类模拟而省略。
另一方面，空间分辨本就约$d_e$，未直接测得完整压力张量；“黏性/轨道效应比惯性重要”来自非绝热区约5 cm大于$d_e$的尺度判断，并非独立完成的欧姆预算。
文章卷期年份是2001，平台2014日期不应改写这一研究的历史位置。
